\documentclass[11pt,letterpaper]{article}
\usepackage[margin=0.60in,top=0.58in,bottom=0.56in]{geometry}
\usepackage{mathptmx}
\usepackage{microtype}
\usepackage{graphicx}
\usepackage{booktabs}
\usepackage{array}
\usepackage{longtable}
\usepackage{amsmath,amssymb}
\usepackage[hidelinks]{hyperref}
\setlength{\parindent}{0pt}
\setlength{\parskip}{2.5pt}
\setlength{\columnsep}{0.25in}
\setcounter{secnumdepth}{0}
\renewcommand{\section}[1]{\vspace{3pt}\noindent\textbf{\underline{#1}}\par\vspace{1pt}}
\renewcommand{\subsection}[1]{\vspace{2pt}\noindent\textbf{#1}\par\vspace{1pt}}
\renewcommand{\familydefault}{\rmdefault}
\providecommand{\tightlist}{\setlength{\itemsep}{0pt}\setlength{\parskip}{0pt}}
\providecommand{\pandocbounded}[1]{\sbox0{#1}\ifdim\wd0>\linewidth\resizebox{\linewidth}{!}{#1}\else\usebox0\fi}
\title{\vspace{-1.2em}\textbf{\MakeUppercase{Neurologically Motivated
Self-Supervised Learning for Bilateral Gait Geometry}}\vspace{-0.5em}}
\author{\normalsize [Authors and affiliations to be inserted in the RISEx template]}
\date{}
\begin{document}
\twocolumn[
  \begin{@twocolumnfalse}
  \maketitle
  \vspace{-1.1em}
  \end{@twocolumnfalse}
]
\subsection{INTRODUCTION}\label{introduction}

Robotic systems that perceive people need motion features whose useful
geometric content can be checked. We test a skeleton JEPA using
bilateral gait. Under reflection, a skeleton's left and right landmarks
exchange places, and a left-minus-right speed contrast changes sign.
This property provides a direct target for evaluating a representation.
Gait asymmetry motivates the setting because it occurs in several health
conditions {[}1{]}, but our score is derived from estimated pose
coordinates and is not a clinical endpoint.

\subsection{MATERIALS AND METHODS}\label{materials-and-methods}

We used 625 GAVD clips from 93 source videos {[}2{]}. Clip-local
preparation keeps all 33 landmarks in a 64-step sequence; four-frame
patches make a 16-by-33 grid. A JEPA predicts hidden teacher features
from visible grid positions {[}3{]}. The bilateral target averages
normalized speed differences for five left--right joint pairs.

The null hypothesis was no gain over an encoder with the same initial
weights. Five outer folds hold out complete source videos from both
encoder training and frozen ridge-readout fitting. Five seeds repeat
each comparison. Matched mask conditions share starts, source draws,
views, and updates. Scores weight source videos equally; paired
resampling uses whole sources.

\subsection{RESULTS AND DISCUSSION}\label{results-and-discussion}

{\def\LTcaptype{none} % do not increment counter
\begin{tabular}{@{}p{0.30\columnwidth}p{0.25\columnwidth}p{0.35\columnwidth}@{}}
\toprule\noalign{}
\begin{minipage}[b]{\linewidth}\raggedright
Check
\end{minipage} & \begin{minipage}[b]{\linewidth}\raggedright
Result
\end{minipage} & \begin{minipage}[b]{\linewidth}\raggedright
Reading
\end{minipage} \\
\midrule\noalign{}
Initial encoder readout & \(R^2=0.223\) & Reference before JEPA \\
Trained teachers & \(R^2=0.101\)--\(0.114\) & Lower in all five arms \\
Trained minus initial & \(-0.109\) to \(-0.122\) & Each reported
interval below zero \\
Correct clip feature target & 375/375 trained checks & Higher
within-clip correspondence \\
Alternative masks & intervals cross zero & No clear gain over random
masks \\
\end{tabular}
}

JEPA training made its predictor discriminate a correct target clip from
a mismatched source more consistently, yet the frozen readout estimated
bilateral movement less accurately. The feature test can use pose, view,
or missingness cues, so it is a diagnostic rather than proof of useful
motion representation. The movement result also depends on the expanded
readout and target preparation; it identifies a limitation of this
pipeline.

\subsection{CONCLUSIONS}\label{conclusions}

Within-clip latent prediction is insufficient evidence that an
articulated-motion representation preserves a downstream geometric
quantity. The proposed evaluation pattern can guide human-aware
perception research before features are transferred to robot-facing
tasks. It does not demonstrate robot performance, control, or clinical
use.

\subsection{REFERENCES}\label{references}

{[}1{]} Patterson KK et al.~\emph{Arch Phys Med Rehabil} 89:304--310,
2008. {[}2{]} Ranjan R et al.~\emph{IEEE Access} 13:45321--45339, 2025.
{[}3{]} Assran M et al.~arXiv:2301.08243, 2023.
\end{document}
